\قسمت{نمایش در \متن‌لاتین{CSS}}

\پاراگراف{}
در رابطه با صفت \متن‌لاتین{display} در \متن‌لاتین{CSS} مقادیر \متن‌لاتین{inline}، \متن‌لاتین{block} و \متن‌لاتین{inline-block} را توضیح دهید.

\پاراگراف{}
\نمره{۲}

% پاسخ پیشنهادی است و پیش از استفاده در تصحیح نیاز به بازبینی دارد.
\begin{پاسخ}

\شروع{فقرات}
\فقره \متن‌سیاه{\متن‌لاتین{inline}}: در جریان متن قرار می‌گیرد، \متن‌لاتین{width} و \متن‌لاتین{height} روی آن اثری ندارد و حاشیه‌های عمودی فضای اضافه ایجاد نمی‌کنند.
\فقره \متن‌سیاه{\متن‌لاتین{block}}: یک خط تازه می‌گیرد، تمام عرض در دسترس را اشغال می‌کند و ابعاد و حاشیه‌ها روی آن اثر دارند.
\فقره \متن‌سیاه{\متن‌لاتین{inline-block}}: مانند \متن‌لاتین{inline} در جریان متن می‌ماند اما مانند \متن‌لاتین{block} ابعاد و حاشیه‌های عمودی می‌پذیرد.
\پایان{فقرات}

هر مورد یک سوم نمره را دارد.

\end{پاسخ}
